\section{A soft vector completion with a bounded coarse remainder}
\label{svc-section}

The exact identities below use the actual mean-zero determinant.
The large-$K$ estimate concerns its first two affine background jets.
Its auxiliary off-shell family retains the effective action $S_M$;
it is not replaced by a bare Wilson Hessian.

\subsection{Exact finite factorization}
\label{svc-factor-section}

Use current counting coordinates, $K=\ell^k$, $\mu=K^{-1}$, and
\begin{equation}
 J=K^2B^T,\qquad J^TJ=I,\qquad \Pi=JJ^T,\qquad P=I-\Pi.
 \label{svc-projectors}
\end{equation}
The actual covariant pair and residual Hessian are
\begin{equation}
 \widehat G=PD^TDP+\mu^2\Pi,\qquad
 Q_B=\mathcal H+DPD^T+a\bar C^T\bar C,
 \qquad \mathcal H=L-C^TqC.
 \label{svc-actual}
\end{equation}
Here $L=S_M''-\lambda\cdot f''$, $q=S_{MK}''$, all Hessians are
left-relative, and $a\ge R_q+1$ is fixed. The actual correction is
$\mathcal D=\mathcal N_0(\log\det\widehat G-\tfrac12\log\det Q_B)$.
Define its soft vector completion and coarse Schur factor by
\begin{equation}
 V=\mathcal H+DD^T+a\bar C^T\bar C
        =Q_B+(DJ)(DJ)^T,
 \qquad S=I-J^TD^TV^{-1}DJ.
 \label{svc-soft-vector}
\end{equation}
The determinant lemma and Woodbury identity give
\begin{equation}
 \det Q_B=\det V\det S,\qquad
 S=(I+J^TD^TQ_B^{-1}DJ)^{-1}.
 \label{svc-schur}
\end{equation}
These are identities between positive matrices on the real small germ,
with normalized holomorphic continuations through flatness.

For the soft scalar matrices
\begin{equation}
 T_s=D^TD+\mu^2\Pi,\qquad M_s=\mu^2J^TT_s^{-1}J,
 \label{svc-soft-scalar}
\end{equation}
the scalar Schur identity is $\det\widehat G=\det T_s\det M_s$.
Thus the exact actual decomposition is
\begin{equation}
 \boxed{\displaystyle
 \mathcal D=\mathcal N_0\left[
 \log\det T_s-\tfrac12\log\det V
                 +\log\det M_s-\tfrac12\log\det S\right].}
 \label{svc-four-factor}
\end{equation}
No massless based scalar factor or omitted Haar normalization occurs.
The matrix $S$ has one fixed color fiber per final scalar site and is
covariant in that coarse frame.

\subsection{Uniform flat spectra and energy resolvents}
\label{svc-flat-section}

A subscript $0$ denotes flat background. The proved mean-zero estimates
are $\Delta_1\le C_\Delta Q_{B,0}$ and
$Q_{B,0}\ge c\mu^2I$. Since
$D_0\Pi_0D_0^T\le D_0D_0^T\le\Delta_1$,
\begin{equation}
 Q_{B,0}\le V_0\le(1+C_\Delta)Q_{B,0},\qquad
 (1+C_\Delta)^{-1}I\le S_0\le I.
 \label{svc-flat-comparison}
\end{equation}
Indeed $\|Q_{B,0}^{-1/2}D_0J_0\|^2\le C_\Delta$ and
\eqref{svc-schur} give the second assertion. The number of coarse
rows has not been used as an estimate for this condition number.

At zero, $qD_c=0$, so the ordinary and corrected constraint terms
agree inside $C^TqC$. With $A_K=\bar C(0)$,
\begin{equation}
 V_0=q_M+D_0D_0^T+A_K^T(aI-q_{MK})A_K
       \ge\Delta_1+A_K^TA_K.
 \label{svc-flat-vector}
\end{equation}
The uniform analytic symbol of $q_M$ vanishes to order two at zero.
Its uniform second moments therefore give $q_M\le C\Delta_1$.
Also $\|A_K\|\le\mu$ and $\|q_{MK}\|\le R_q$.
Combining these facts with the gap gives
\begin{equation}
 c(\Delta_1+\mu^2I)\le V_0\le C(\Delta_1+\mu^2I),
 \label{svc-energy-comparison}
\end{equation}
on the chosen infinite structural system and its explicit periodizations. The finite actual lower bound, gap and heat estimate below remain valid without identifying arbitrary finite winding kernels with those periodizations.

For the spatial argument, $q_M$ is the kernel of the chosen infinite structural system. The same conclusions hold for its explicitly periodized kernels; they are not asserted for unrelated finite winding kernels merely by sampling. For later source estimates the spatial exponent must also be uniform
in final units. If $\delta(p)$ is the scalar lattice symbol, the common
analytic strip and double zero of $q_M$ imply
\begin{equation}
 \|q_M(p+i\mu\theta)-q_M(p)\|
 \le C\{\mu|\theta|\sqrt{\delta(p)}+\mu^2|\theta|^2\}.
 \label{svc-fine-twist}
\end{equation}
Near zero this is Taylor's formula; away from zero it follows from
bounded derivatives. The ordinary longitudinal term has the same bound.
Conjugating fine fields by $e^{\theta\cdot x/K}$ and coarse fields by
$e^{\theta\cdot y}$, straight-average row and column counting gives
\begin{equation}
 \|A_{K,\theta}\|\le C\mu,\qquad
 \|A_{K,\theta}-A_K\|\le C|\theta|\mu.
 \label{svc-average-twist}
\end{equation}
The coarse symbol $q_{MK}$ has a uniform strip. Its coarse correction
therefore changes by $O(|\theta|\mu^2)$.
Measured between two $(\Delta_1+\mu^2I)^{-1/2}$ factors, the total
change is $O(|\theta|)$. Equations \eqref{svc-energy-comparison}--
\eqref{svc-average-twist} give a common physical strip on which all
energy sandwiches of $V_0^{-1}$ are bounded uniformly. The energy legs
are the ordinary componentwise first differences and $\mu I$.

Bloch decomposition by $K$-blocks turns these into bounded
operator-valued analytic functions. Fourier integration gives uniform
exponentially summable final-block kernels. The internal fiber may
have dimension $K^4$, but the bounds are Hilbert operator norms; no
entrywise sum over that fiber is taken. Explicit periodizations obey the same bounds. The established fixed-$K$ local-limit interface identifies the resulting infinite source coefficient with the actual local Maxwell coefficient; it does not identify arbitrary finite kernels with these periodizations. In particular $D_0^TV_0^{-1}D_0$, $S_0$ and, after shrinking
the exponent using \eqref{svc-flat-comparison}, $S_0^{-1}$ have uniform
weighted block bounds.

Min--max also gives the flat heat estimate
\begin{equation}
 N^{-4}\operatorname{Tr}e^{-uK^2V_0}
 \le C\left\{\frac{K^4}{(1+K^2u)^2}+N^{-4}\right\}e^{-cu}.
 \label{svc-flat-heat}
\end{equation}
It follows from $V_0\ge Q_{B,0}$ and the mean-zero heat theorem.
The same Mellin integration gives bounded inverse moments for powers
$p>2$, the logarithmic bound at $p=2$, and exponential final-time tails.
These flat estimates alone say nothing about a magnetic derivative.

\subsection{The auxiliary family and the exact affine matching}
\label{svc-auxiliary-section}

For bounded physical $C^1$ potentials $A$ put $U_A=e^{i\mu A}$ and
\begin{align}
 \widetilde V(A)
  &=S_M''(U_A)+D_AD_A^T
       +\bar C_A^T\{aI-S_{MK}''(M_K(U_A))\}\bar C_A,
       \label{svc-auxiliary-vector}\\
 \widetilde S(A)
  &=I-J_A^TD_A^T\widetilde V(A)^{-1}D_AJ_A.
       \label{svc-auxiliary-coarse}
\end{align}
Every Hessian in these formulas is the actual relative-coordinate
classical Hessian. The family is covariant, but it is not declared equal
to \eqref{svc-soft-vector} off shell.

On the established commuting affine probe, the actual fine and coarse
currents and the multiplier vanish through the two jets used in the
Maxwell extraction. Thus $qD_c=0$ and
\begin{equation}
 C^TqC=\bar C^Tq\bar C,\qquad L=S_M''
 \quad\hbox{through the first two affine jets}.
 \label{svc-on-shell-match}
\end{equation}
The actual smooth stationary polynomial jets are the corresponding
fine affine connection. Hence $V,S$ and
$\widetilde V,\widetilde S$ have identical first two affine jets.
This is the precise interface used below.

\subsection{Relative vertices with localized sources}
\label{svc-vertices-section}

Use complete physical source fields with
\[
 a_* =\|A\|_\infty+
       \mu^{-1}\max_{\text{nearest, same component}}|A(e')-A(e)|.
\]
The corrected relative-Hessian theorem supplies, for $j=1,2$,
\begin{align}
 [t^j]S_M''(e^{it\mu A})
     &=d_0^TZ_j+Z_j^Td_0+E_j,\notag\\
 \|Z_1\|&\le C\mu a_*,&\quad \|E_1\|&\le C\mu^2a_*,\notag\\
 \|Z_2\|&\le C\mu^2a_*^2,&\quad \|E_2\|&\le C\mu^2a_*^2.
 \label{svc-relative-vertices}
\end{align}
Here $d_0$ is the componentwise gradient of a bond field. These are
source-rooted block estimates, including physical weights.
Their proof uses the ordered Laurent factorization
$q_M=d_0^TH_0d_0$ and a symmetric covariant completion. A common
componentwise-difference formula would be false for the shifted
component-skew kernel and is not used.

The first two $D_A$ jets are multiplication operators of sizes
$O(\mu a_*)$ and $O(\mu^2a_*^2)$. Their contribution to
$D_AD_A^T$ has the same form as \eqref{svc-relative-vertices}.
The actual moving-reference products obey
$\|\bar C_A\|=O(\mu)$ with first two physical-background derivatives
of the same scale. Their rows and sources lie in two final blocks.
The ordinary output $M_K(e^{i\mu A})$ is uniformly analytic and
$O(\|A\|)$ in final units, with the same fixed physical range.
The previously proved relative coarse Hessian kernels are uniform
in $MK$. Therefore the last term of
\eqref{svc-auxiliary-vector}, including its first two derivatives,
is a zero-order $O(\mu^2)$ vertex with exponentially rooted sources.
These conclusions use corrected moving tangents, not an estimate
for the full uncorrected ordinary adjoint.

The normalized block map $J_A$ and its first two physical derivatives
have uniform own-block operator bounds. Keep the factorization
\begin{equation}
 R(A)=D_A^T\widetilde V(A)^{-1}D_A,
 \qquad \widetilde S(A)=I-J_A^TR(A)J_A.
 \label{svc-avoid-interface}
\end{equation}
The derivatives of $J$ stay outside $R$. In particular, no raw derivative
of $DJ$ at a sharp block interface is estimated before using a resolvent.
Since $V_0\ge D_0D_0^T$ and $V_0\ge c\mu^2I$,
\begin{equation}
 \|V_0^{-1/2}D_0\|\le1,\qquad
 \|V_0^{-1/2}D'_0\|\le Ca_*.
 \label{svc-gauge-legs}
\end{equation}
For example,
\begin{equation}
 R'=D'^TV_0^{-1}D_0+D_0^TV_0^{-1}D'
           -D_0^TV_0^{-1}V'V_0^{-1}D_0.
 \label{svc-first-response}
\end{equation}
Differentiating once more yields the ordered two-$V'$ terms, one
$V''$ term, and the $D''$ and mixed $D',V'$ terms. Insert
\eqref{svc-relative-vertices} between the uniform weighted energy
resolvents. Each term is bounded with its physical source norms,
by \eqref{svc-gauge-legs}; this proves both first and second response
bounds without a fine-coordinate dimension factor.

For unit directions supported in smooth physical collars $x,y$, a
spare exponent attaches the sources to their resolvent legs and gives
\begin{align}
 \|\widetilde S'_x(i,j)\|&\le Ce^{-m\tau(i,j,x)},\notag\\
 \|\widetilde S''_{x,y}(i,j)\|
      &\le Ce^{-m\tau(i,j,x,y)}.
 \label{svc-rooted-sources}
\end{align}
Here $\tau$ is connecting-tree length including every displayed operator
endpoint and source. With a smaller exponent, these imply
$\sum_{i,j}e^{m'[d(i,x)+d(j,x)]}\|\widetilde S'_x(i,j)\|\le C$
and $\sum_{i,j}\|\widetilde S''_{x,y}(i,j)\|\le Ce^{-m'd(x,y)}$.
The full rooted estimate, rather than only these summed consequences,
is retained in the diagonal-density argument below.
The indices are final-block labels, with fixed color fibers.
The same source bounds hold after multiplying by $J,J',J''$, because
these are normalized own-block operators. Completeness of the fine
field collars is essential; sharp background cutoffs are not differentiated.

\subsection{The coarse logarithm and its actual Maxwell bound}
\label{svc-maxwell-section}

Equations \eqref{svc-flat-comparison} and
\eqref{svc-rooted-sources} give for the total Hessian
\begin{equation}
 K_S(x,y)=\operatorname{Tr}\left[
 S_0^{-1}\widetilde S''_{x,y}
 -S_0^{-1}\widetilde S'_xS_0^{-1}\widetilde S'_y\right],
 \qquad |K_S(x,y)|\le Ce^{-m' d(x,y)}.
 \label{svc-log-hessian}
\end{equation}
Both traces are coarse traces with fixed color dimension.

For the local diagonal density one must not cyclically rearrange a
partial trace. Put $R_t=(S_0+tI)^{-1}$. Its second source derivative
at root zero is the partial trace of the integral over $t\ge0$ of
\begin{equation}
 R_t\widetilde S''_{x,y}R_t
 -R_t\widetilde S'_xR_t\widetilde S'_yR_t
 -R_t\widetilde S'_yR_t\widetilde S'_xR_t.
 \label{svc-density-resolvent}
\end{equation}
The weighted bound is $\|R_t\|_m\le C/(1+t)$.
Thus its rooted density kernel is bounded by
$Ce^{-m''\tau(0,x,y)}$. This gives uniform polynomial source moments,
not merely decay in the source separation.

Covariance makes each diagonal color trace invariant. Apply the
rooted bounds to a smooth physical partition of an affine potential;
the collar norms grow only linearly with their distance from zero and
are summable against the rooted exponential. The density Ward and
reanchoring identity then yields
\begin{equation}
 \left|\beta\left(\tfrac12\mathcal N_0\log\det S\right)\right|
       \le C_S,
 \label{svc-bounded-remainder}
\end{equation}
uniformly in $M,K$, in the actual joint infinite structural systems
and polynomial-jet convention. The transfer from $\widetilde S$ to the
actual $S$ uses exactly \eqref{svc-on-shell-match}. The linear color jet
vanishes by semisimple conjugation invariance.

Combining \eqref{svc-four-factor} with the scalar coarse theorem gives
\begin{equation}
 \beta(\mathcal D)
  =\beta\!\left(\mathcal N_0
       [\log\det T_s-\tfrac12\log\det V]\right)+O(1),
 \label{svc-final-reduction}
\end{equation}
with a remainder uniform in old depth and endpoint. For the vector
coefficient, the auxiliary expression
\eqref{svc-auxiliary-vector} is valid through its two affine jets.
The fine soft pair has not yet been compared to a standard massive
vector/ghost reference. A fixed-integrand universality theorem alone
does not supply bounds uniform in this family. This module proves
neither that remaining logarithmic comparison nor its sign, and it
does not assert a full $K$-independent nonlinear background domain.
Finite-volume winding and normalized-extraction issues retain their
separate regulator hypotheses.

